\documentclass[11pt]{article}
\usepackage[T1]{fontenc}
\usepackage{lmodern}
\usepackage{microtype}
\usepackage{amsmath,amssymb}
\usepackage{booktabs,longtable}
\usepackage[round,authoryear]{natbib}
\usepackage{array}
\usepackage{xurl}
\usepackage[margin=0.88in]{geometry}
\usepackage{enumitem}
\usepackage{graphicx}
\usepackage[hidelinks]{hyperref}
\setlist{nosep}
\setlength{\parskip}{0.24em}
\setlength{\parindent}{1.2em}
\newcommand{\LCDM}{$\Lambda$CDM}
\newcommand{\methodsdoi}{10.5281/zenodo.22022089}
\newcommand{\draftfivedoi}{10.5281/zenodo.22025329}
\newcommand{\draftsixdoi}{10.5281/zenodo.22031121}
\newcommand{\vsevendoi}{10.5281/zenodo.22031627}
\newcommand{\veightdoi}{10.5281/zenodo.22032094}
\newcommand{\tauinf}{\tau_{\infty,B}}
\newcommand{\kreact}{\kappa_{\rm r}}

\hypersetup{
 pdftitle={A Static-Universe Two-Channel Photon Propagation Model Confronted with DES-SN5YR - Version 8},
 pdfauthor={T. McSheery},
 pdfsubject={DES-SN5YR, static Euclidean geometry, redshift, time dilation, chromatic transfer, Tolman surface brightness, energy closure},
 pdfkeywords={DES-SN5YR, supernovae, static universe, photon scattering, time dilation, Tolman test, falsification, CMB, photon counting}
}

\title{A Static-Universe Two-Channel Photon Propagation Model Confronted with DES-SN5YR\\
\large A Standalone Adversarial Test of Redshift, Temporal Remapping, Chromatic Transfer, Surface Brightness, and Energy Closure}
\author{T.~McSheery\\
\small PhaseSpace Inc., San Leandro, CA, USA\\
\small \href{mailto:TracyMcSheery@gmail.com}{TracyMcSheery@gmail.com}}
\date{20 August 2026}

\begin{document}
\maketitle

\begin{center}
\small
Version 8.0 --- human-circulation freeze\\
Current archive DOI: \href{https://doi.org/\veightdoi}{\veightdoi}\\
Prior records preserved as provenance: V7 \href{https://doi.org/\vsevendoi}{\vsevendoi}; Draft 6 \href{https://doi.org/\draftsixdoi}{\draftsixdoi}; Draft 5 \href{https://doi.org/\draftfivedoi}{\draftfivedoi}\\
Independent methods archive: \href{https://doi.org/\methodsdoi}{\methodsdoi}\\
Canonical code and audit trail: \href{https://github.com/tracyphasespace/Static-Universe-DES-SN5YR}{github.com/tracyphasespace/Static-Universe-DES-SN5YR}
\end{center}

\begin{abstract}
This paper asks a deliberately adversarial question: can the principal Type Ia supernova observables normally interpreted as evidence for metric expansion be reproduced by a non-expanding Euclidean spatial geometry with two photon--vacuum propagation channels, while surviving the classical tests that eliminate simpler tired-light models? ``Static'' here means no large-scale metric expansion; Version 8 explicitly allows source populations and cosmic contents to evolve with emission epoch and therefore does not assume a stationary eternal galaxy ensemble.

The first propagation channel is coherent and direction preserving. At effective level it has two observable faces that must ultimately be derived from one microscopic vertex. An inelastic frequency-shift face gives $d\ln\nu/dD=-K/c$, hence $D=(c/K)\ln(1+z)$ for steady and transient sources. A distinct reactive group-delay face is represented as a direction-resolved phase-space susceptibility $\chi(\mathbf{x},\hat{\mathbf{k}},t)$ responding to temporal structure above the saturated DC background of a propagation cell. Along one supernova ray its characteristic group-slowness factor is $n_g=1+\kreact K\xi$, with $\xi=t-r/c$ measured behind the transient front. The leading front remains on $n_g=1$ while a trailing independent photon obeys $d\ln\Delta t/dD=\kreact K/c$, giving $\Delta t_{\rm obs}=\Delta t_{\rm em}(1+z)^{\kreact}$. DES-SN5YR measures $b=1.003\pm0.005\,(\mathrm{stat})\pm0.010\,(\mathrm{sys})$, requiring $\kreact\simeq1$. The kinematics are therefore explicit and causal without optical coherence between photons; the unresolved microscopic problem is to derive the equality between the reactive rate and the inelastic drag rate from the vacuum action rather than impose it phenomenologically.

A second, non-forward channel has $\sigma_{\rm nf}\propto E^{1/2}\propto\lambda^{-1/2}$. The physically defined rest-$B$ asymptotic optical depth corresponding to the DES fit is $\tauinf=0.594$, removing 10.3\% of the coherent beam at $z=0.5$ and 16.0\% at $z=1$. On the released DES-SN5YR Hubble diagram, the resulting one-parameter static distance law is nearly degenerate with one-parameter flat \LCDM: $\chi^2\simeq1641.1$ versus 1640.1 on the default 1,768-object diagnostic sample, $\Delta\chi^2=+1.1$, with both $\chi^2/\mathrm{dof}\simeq0.93$.

Several tempting supporting mechanisms were then tested rather than retained. A full Hsiao-SED $\times$ measured-DECam-filter calculation shows that Wien/passband slicing plus the chromatic opacity changes fitted stretch by at most $5\times10^{-4}$ and peak epochs by at most 0.05 d when the temporal remap is disabled; it cannot manufacture the observed $b\simeq1$. A preregistered iron-group secondary-maximum clock measured on raw DES $i/z$ photometry gives $b_2=0.52\pm0.78$ from 120 detections: explicitly inconclusive, consistent with both universal remapping and no dilation because DES cadence and the $z\lesssim0.3$ visibility window are underpowered.

The strongest present empirical threat is the Tolman surface-brightness test. The model predicts $D_A=D=(c/K)\ln(1+z)$ and $I\propto(1+z)^{-2}e^{-\tau_\lambda}$. Re-expressing the published Lubin--Sandage distance sensitivity in the static frame removes the nominal $I$-band crisis but moves tension to $R$; the approximately frame-invariant observed spread $n_I-n_R=0.78\pm0.21$ differs from the model prediction $\simeq-0.02$ by about 3.7$\sigma$ before a full static-frame re-reduction. The paper also closes the previously omitted forward-energy ledger: at $z=1$ the coherent frequency shift transfers 50\% of each photon's emitted energy out of the observed photon, while the non-forward branch removes a further 16\% of the surviving beam in rest $B$. Applying the same frequency drift to the microwave bath predicts the independently observed scaling $T_{\rm CMB}(z)=T_0(1+z)$; at $z=2.41837$ it gives 9.315 K versus $9.15\pm0.72$ K from CO excitation. This is a consistency success shared with FLRW, not a discriminator. Preserving Planck normalization in a static volume then requires an additional shape-preserving photon-number relaxation at rate $3K$, whose microphysics is not identified here, while the CMB acoustic/polarization/lensing and BAO-correlated structure remains unmodeled. Version 8 therefore does not claim a completed cosmology. It claims a sufficiently specified supernova alternative with explicit adverse tests and a finite microscopic closure problem suitable for human review.
\end{abstract}

\section{Thesis, scope, and terminology}
The thesis tested in this paper is:
\begin{quote}
\textbf{Thesis.} A non-expanding Euclidean spatial geometry with two photon--vacuum propagation channels can reproduce the observed achromatic redshift, the approximately $(1+z)$ temporal broadening of supernova signals, and the DES-SN5YR luminosity-distance relation without a dark-energy acceleration term. The same transfer law must then survive independent tests of chromaticity, angular distance, surface brightness, radiance, image sharpness, source-population evolution, and energy conservation.
\end{quote}

Three distinctions are load bearing.

First, \emph{static geometry is not stationary contents}. Light still has finite travel time. For the distance law derived below,
\begin{equation}
 t_{\rm lb}(z)=\frac{D(z)}{c}=K^{-1}\ln(1+z).
 \label{eq:lookback}
\end{equation}
Version 8 therefore adopts the logically available branch ``static spatial geometry, evolving contents.'' Observed stellar-population evolution, the cosmic star-formation history, quasar evolution, and related redshift trends are not switched off; they become constraints on the history of matter and sources evaluated against Eq.~\eqref{eq:lookback}. This costs any claim that the observed universe is a strictly stationary eternal ensemble.

Second, the model is not classical tired light. Classical energy-loss models generally predict no macroscopic time dilation and often introduce chromaticity or image blur. Here redshift, transient timing, and non-forward chromatic depletion are separate observables with separate failure modes.

Third, the paper follows an anti-confirmation-bias rule. A mechanism that fails a registered test is retained as a null result, not reinterpreted. An underpowered measurement is labeled inconclusive. An adverse observable is placed in the same status ledger as favorable ones. The public repository preserves preregistration amendments, implementation defects, and superseded mechanisms.

\section{Physical motivation without microscopic overclaiming}
Nonlinear electromagnetic propagation is not a forbidden category. The Euler--Heisenberg effective action predicts nonlinear QED interactions, and SLAC E-144 observed nonlinear Compton processes and multiphoton pair production using an intense laser and a high-energy electron beam \citep{Burke1997,Bamber1999}. Those experiments do not establish the cosmological interaction proposed here, and they do not supply its rate. They establish only that a nonlinear vacuum electromagnetic response is physically meaningful.

The present object is therefore an \emph{effective propagation model}. Its coefficients are constrained by astronomical observations. A successful microscopic theory must derive them from the underlying vacuum dynamics while respecting spectral-line coherence, image sharpness, polarization, local laboratory constraints, energy conservation, and the CMB. Where that derivation is absent, the paper says so explicitly.

\section{Effective two-face forward vertex}
\subsection{Inelastic coherent face: universal frequency drift}
The coherent forward interaction is assumed direction preserving and achromatic in fractional energy loss. At the effective level the same drift law applies to directed source photons and to photons in the diffuse bath,
\begin{equation}
 \frac{dE}{dD}=-\frac{K}{c}E,
 \qquad
 \frac{d\nu}{dD}=-\frac{K}{c}\nu.
 \label{eq:drag}
\end{equation}
Integrating gives
\begin{equation}
 1+z=\frac{\nu_{\rm em}}{\nu_{\rm obs}}=e^{KD/c},
 \qquad
 \boxed{D(z)=\frac{c}{K}\ln(1+z)}.
 \label{eq:distance}
\end{equation}
The same fractional shift applies to all spectral coordinates,
\begin{equation}
 \lambda_{\rm obs}=(1+z)\lambda_{\rm em}.
\end{equation}
A thermal spectral shape therefore has its Wien peak mapped to the coordinate temperature $T/(1+z)$. This statement concerns spectral coordinates only; full Planck radiance normalization is a separate test in Section~\ref{sec:tolman}.

Equation~\eqref{eq:drag} is the redshift face of the effective vertex. It does not require a preconditioned corridor and remains active for steady galaxy light. It also creates a substantial energy ledger: by redshift $z$, the photon has transferred the fraction
\begin{equation}
 f_{\rm fwd}(z)=1-\frac{1}{1+z}=\frac{z}{1+z}
 \label{eq:fwdloss}
\end{equation}
out of the observed photon. At $z=1$, $f_{\rm fwd}=0.5$. Section~\ref{sec:energy} treats the destination of this energy as a required closure condition.

\subsection{Reactive face: direction-resolved phase-space memory}
The temporal problem is different. A time-translation-invariant, unchanged medium cannot stretch the arrival interval between mutually incoherent photons emitted days apart. The model therefore assigns a second, reactive face to the same effective interaction: a weak, \emph{direction-resolved} group-delay susceptibility produced by temporal structure in the radiation field above its saturated DC background.

The memory variable is not a bulk scalar refractive index. It is a phase-space response $\chi(\mathbf{x},\hat{\mathbf{k}},t)$ associated with a propagation-direction cell (and, microscopically, potentially polarization). Foreign rays crossing the same spatial point occupy different $\hat{\mathbf{k}}$ cells and therefore do not automatically add to the same memory state. Along the one-ray characteristic of a supernova, define
\begin{equation}
 \xi=t-\frac{r}{c}
 \label{eq:xi}
\end{equation}
and the effective group slowness factor
\begin{equation}
 n_g(\mathbf{x},\hat{\mathbf{k}},t)\big|_{\rm ray}
 =1+\kreact K\xi,
 \qquad 0<\xi<T_{\rm sat}.
 \label{eq:ng}
\end{equation}
Here $n_g$ is shorthand for the one-ray characteristic of the reactive susceptibility, not a scalar property of the bulk vacuum. The coefficient $\kreact$ measures the reactive rate relative to the inelastic drift rate $K$.

The leading transient front sits at $r=ct$, hence $\xi=0$ and $n_g=1$ for its entire trajectory. Explicitly,
\begin{equation}
 \frac{dn_g}{dt}=\partial_t n_g+\dot r\,\partial_r n_g
 =\kreact K+c\left(-\frac{\kreact K}{c}\right)=0.
 \label{eq:totalderiv}
\end{equation}
This is the point that resolves the earlier distance-law objection: the local transient index excursion can be tiny even though the propagation effect accumulates over a cosmological path. The leading front is never slowed to $c/(1+z)$; it rides the moving $\xi=0$ surface.

Now consider a trailing photon emitted independently a time $\Delta t$ later. It samples $\xi=\Delta t$ and therefore, to first order,
\begin{equation}
 \frac{d\Delta t}{dD}=\frac{\kreact K}{c}\Delta t.
\end{equation}
Thus
\begin{equation}
 \boxed{\Delta t_{\rm obs}=\Delta t_{\rm em}e^{\kreact KD/c}
       =\Delta t_{\rm em}(1+z)^{\kreact}.}
 \label{eq:timemap}
\end{equation}
The measured supernova time-dilation exponent is therefore
\begin{equation}
 b=\kreact.
 \label{eq:bkr}
\end{equation}
White et al. measure \citep{White2024}
\begin{equation}
 b=1.003\pm0.005\ ({\rm stat})\pm0.010\ ({\rm sys}),
 \label{eq:white}
\end{equation}
so the effective response requires $\kreact\simeq1$.

This formulation uses no mutual optical coherence between photons. The corridor field is the shared memory variable. It is also more honest than the earlier claim that $b=1$ was already a microscopic identity: \emph{at effective level} the measured coefficient fixes the equality of the reactive and inelastic rates. Deriving $\kreact=1$ from a single microscopic vacuum vertex is still owed.

\subsection{What defines $\xi=0$ in a permanently illuminated universe?}
A sharpened referee objection asks what ``the transient front'' means when the line of sight has been illuminated by the host galaxy and diffuse backgrounds for gigayears. The effective answer is a two-level response:

\begin{enumerate}[label=(\roman*)]
\item The inelastic coherent face, Eq.~\eqref{eq:drag}, is the universal frequency-drift channel. It redshifts steady galaxy light without requiring growth of the reactive memory.
\item The reactive face is an AC response in a direction-resolved phase-space cell. The steady host contribution defines the saturated DC background of that cell; variability $\delta f_{\hat{k}}(t)$ activates the time-dependent response without summing unrelated transients from other directions into the same scalar field. For a supernova, $\xi=0$ is the onset of the event's directional transient above its local DC state.
\end{enumerate}

A generic effective kinetics, not yet a Lagrangian derivation, is therefore written
\begin{equation}
 \partial_t\chi_{\hat{k}}
 =\mathcal{R}[\delta f_{\hat{k}}]
 -\frac{\chi_{\hat{k}}}{T_{\rm relax}},
 \qquad n_g\big|_{\hat{k}}=1+\frac{\chi_{\hat{k}}}{2}.
 \label{eq:kinetic}
\end{equation}
A previously used Heaviside/saturation switch is only an existence sketch; it is not retained as the physical response law because a real response must handle both rising and declining light-curve structure. The operator $\mathcal{R}$ must select the appropriate direction/polarization phase-space cell while keeping the activated-cell rate set by the medium rather than by beam intensity. Equation~\eqref{eq:kinetic} specifies what the microphysics must produce, not evidence that it does. A successful underlying theory must derive the universal drift rate $K$, the empirical reactive coefficient $\kreact\simeq1$, the angular and temporal coarse-graining scale, saturation behavior, and relaxation time from one vacuum response.

This sourcing split removes the false dichotomy between a global $n(t)$ and independent per-wavepacket wakes. Redshift is a universal frequency-drift process, while timing is a direction-resolved transport memory shared by later photons in the same event cell. Whether one microscopic vertex can actually generate both faces with the required coefficient equality is the central theoretical debt of Version 8.

\section{Non-forward chromatic channel}
The second propagation channel removes photons from the coherent image-forming beam with an energy-dependent probability. The working scaling is
\begin{equation}
 \sigma_{\rm nf}(E)\propto E^{1/2}\propto\lambda^{-1/2}.
 \label{eq:sigma}
\end{equation}
Along the redshifting path this gives the rest-frame optical depth
\begin{equation}
 \tau_\lambda(z)=\tau_{\infty,\lambda}
 \left[1-(1+z)^{-1/2}\right],
 \qquad
 \tau_{\infty,\lambda}=\tauinf\left(\frac{440\,{\rm nm}}{\lambda}\right)^{1/2}.
 \label{eq:tau}
\end{equation}
The physical magnitude attenuation is
\begin{equation}
 \Delta m_\lambda=\frac{2.5}{\ln 10}\tau_\lambda.
 \label{eq:magop}
\end{equation}
Earlier versions used a $5/\ln10$ bookkeeping coefficient; the fitted amplitude is unchanged, but its physical interpretation doubles. The rest-$B$ asymptotic optical depth is therefore
\begin{equation}
 \boxed{\tauinf=0.594}.
\end{equation}
It removes 10.3\% of rest-$B$ coherent-beam photons at $z=0.5$ and 16.0\% at $z=1$.

The complete observer-frame transfer operator used in the synthetic tests is
\begin{equation}
 F_{\rm obs}(\lambda_o,t_o)=A(z)
 e^{-\tau(\lambda_o/(1+z),z)}
 S\!\left(\frac{\lambda_o}{1+z},\frac{t_o}{(1+z)^{\kreact}}\right),
 \label{eq:transfer}
\end{equation}
where $S(\lambda,t)$ is the intrinsic SN spectral-time surface. No factor in Eq.~\eqref{eq:transfer} knows which atomic species produced a feature; species physics determines $S$, while propagation acts on the entire surface. That universality is falsifiable.

\section{A mechanism tested and killed: passbands do not fake $b=1$}
A tempting explanation was that Wien displacement, finite observer filters, and phase-dependent chromatic attenuation might make an intrinsically undilated transient appear broader. This was tested directly rather than retained.

Using the Hsiao SN-Ia spectral time series \citep{Hsiao2007}, measured DECam $g,r,i,z$ throughput curves, Eq.~\eqref{eq:tau}, and the spectral coordinate shift, we generated synthetic light curves with the temporal remap in Eq.~\eqref{eq:transfer} disabled. The recovered stretch changes by at most
\begin{equation}
 |\Delta s|\le5\times10^{-4}
\end{equation}
for $z\le0.8$ in every band, and recovered peak epochs move by at most 0.05 d. An independently preregistered inverse-injection test likewise returned $\hat b=0$ for time-independent opacity, even when that opacity was amplified far beyond the fitted value. The structural reason is simple: a time-independent multiplicative spectral weight can reshape color and amplitude but does not remap the temporal coordinate.

With the temporal remap enabled at $\kreact=1$, the same transfer calculation recovers $(1+z)$ stretch to within 0.05\% in each DES band. The chromatic branch therefore survives with a smaller, cleaner job: amplitude and color transfer, not time dilation.

For phase-matched rest-frame windows near 330 and 650 nm the model predicts
\begin{equation}
 Q(z)\equiv\frac{[F_{330}/F_{650}]_z}{[F_{330}/F_{650}]_{z\simeq0}}
 =e^{-\Delta\tau(z)},
\end{equation}
with
\begin{equation}
 Q(0.25)=0.979,\qquad Q(0.5)=0.964,\qquad Q(1)=0.944.
 \label{eq:Q}
\end{equation}
The propagation component is predicted to be phase-flat at fixed rest wavelengths. Intrinsic line formation and recombination evolve strongly with phase; phase-flatness is therefore part of the discriminant, not merely the amplitude trend.

\section{DES-SN5YR distance comparison}
\subsection{Standalone likelihood definition}
The analysis uses the public DES-SN5YR v1.2 Hubble-diagram products, the official $z_{\rm HD}/z_{\rm HEL}$ convention, and the full released STAT+SYS covariance. The absolute-magnitude/distance-scale offset is analytically profiled. The default diagnostic sample excludes entries with effectively zero inverse-variance weight ($\mathrm{MUERR}>10$), leaving $N=1768$; retaining all 1,829 released entries changes no quoted headline number.

For the static law,
\begin{align}
 \mu_{\rm stat}(z_{\rm HD},z_{\rm HEL})
 &=5\log_{10}\!\left[(1+z_{\rm HEL})\ln(1+z_{\rm HD})\right]\\
 &\quad+\frac{2.5}{\ln10}\tauinf
 \left[1-(1+z_{\rm HD})^{-1/2}\right]+\mathcal M.
 \label{eq:mustat}
\end{align}
Flat \LCDM is fitted with one shape parameter $\Omega_m$ and the same profiled offset. The physical opacity prior is non-negative; the numerical fit is bounded well above the optimum.

The released-vector best fits are approximately $\Omega_m=0.352$ for flat \LCDM and $\eta=0.297$ in the historical magnitude-amplitude convention for the static law, equivalent to the physical $\tauinf=0.594$ used here. The resulting fit values are
\begin{equation}
 \chi^2_{\Lambda{\rm CDM}}\simeq1640.1,
 \qquad
 \chi^2_{\rm stat}\simeq1641.1,
\end{equation}
so
\begin{equation}
 \boxed{\Delta\chi^2\equiv\chi^2_{\rm stat}-\chi^2_{\Lambda{\rm CDM}}=+1.1},
 \label{eq:headline}
\end{equation}
with both $\chi^2/\mathrm{dof}\simeq0.93$. The maximum offset-removed separation of the two best-fit curves over the DES redshift range is approximately 0.055 mag.

This is a deliberately narrow result. It is a one-dimensional SN distance comparison, not a global cosmological fit. BAO, CMB anisotropies, lensing, growth, nucleosynthesis, and source-population evolution are separate constraints and are explicitly retained in the open ledger below.

\subsection{Bias-correction sensitivity is methodological, not cosmological evidence}
For completeness, the companion methods analysis reconstructs a diagnostic vector before the event-level DES BBC magnitude correction. The pairwise statistic changes from $+1.1$ on the released vector to $+12.8$ on that diagnostic vector, a total processing leverage of $-11.7$ $\chi^2$ units. The pre-BBC vector contains known survey-selection bias and the released covariance belongs to the corrected product, so this number is not evidence for the static model. It motivates a matched forward/BBC simulation under the alternative law. The full sign audit, covariance projections, robustness probes, and frozen-result CI live in the independent methods archive DOI \methodsdoi; no result in the present paper requires accepting a cosmological interpretation of that leverage.

The strongest existing DES counterargument is its explicit reference-cosmology validation. Camilleri et al. reran the bias-correction machinery under multiple reference cosmologies spanning distance-modulus variations of order 0.15 mag and found recovered cosmological shifts subdominant to the DES statistical uncertainty \citep{Camilleri2024}. The offset-removed separation between the released best-fit static curve and best-fit flat \LCDM is only about 0.055 mag, below that tested amplitude scale. This is relevant evidence that the BBC layer is not arbitrarily fragile. It is not, however, the same experiment as propagating the present non-FLRW transfer law through the simulation and selection pipeline: the published reference set is FLRW-family and the formal validation is shape dependent, not merely an amplitude bound. Version 8 therefore treats the matched alternative-reference rerun as a completion test, not as an assumed failure of DES calibration.

\section{Temporal universality and the executed T4 test}
White et al. find Eq.~\eqref{eq:white} from 1,504 DES Type Ia supernovae \citep{White2024}. Any static model with $b=0$ is therefore excluded. In the present effective model, Eq.~\eqref{eq:timemap} makes this a direct measurement of $\kreact$ rather than an assumed optical-coherence effect.

Blondin et al. measured multi-epoch spectra of high-redshift SNe Ia and found spectral aging consistent with a $1/(1+z)$ rate \citep{Blondin2008}. This is a physically different clock and rejects no-dilation models, but it still does not distinguish FLRW from a universal non-metric remap.

To add a species-sensitive internal clock, we preregistered a test based on the Type Ia $i/z$ secondary maximum, which is associated with iron-group recombination and opacity/flux redistribution rather than one isolated atomic transition \citep{Jack2015}. The frozen primary regression was
\begin{equation}
 y=a+b_2x+\gamma(x_1-\bar x_1)+\epsilon,
 \quad
 y=\ln\!\frac{\Delta t_{\rm obs}}{\Delta t_{\rm model}(z,{\rm band})},
 \quad x=\ln(1+z).
 \label{eq:t4}
\end{equation}
The universal-remap prediction is $b_2=1$ and the no-dilation null is $b_2=0$.

The executed result is
\begin{equation}
 \boxed{b_2=0.52\pm0.78,\qquad N=120.}
 \label{eq:t4result}
\end{equation}
It lies 0.6$\sigma$ from 1 and 0.7$\sigma$ from 0. Neither preregistered kill condition fires. The result is therefore \emph{inconclusive}, not evidence for a mixture and not positive evidence for the model.

The run nevertheless established why DES $griz$ cannot decide this test. Median per-SN secondary-maximum timing uncertainty was 5.3 d rather than the forecast 2 d; intrinsic bump diversity contributed about an 8 d equivalent spread; and the feature leaves DES $griz$ beyond roughly $z=0.3$. Approximately 25 times the effective statistics would be required for a 3$\sigma$ 0-versus-1 distinction under the realized noise. The detection-efficiency difference between low- and high-$z$ bins was only 0.8$\sigma$.

The complete audit trail is public. Pre-data amendments added an intercept, a frozen GP peak finder, censoring, and a hypothesis-neutral search window. The first executable pipeline then exposed single-epoch difference-imaging artifacts; a post-data 5$\sigma$ leave-one-out screen and three implementation repairs were disclosed before any valid $b_2$ result existed. Version 8 keeps that history because the ability to report an unhelpful result is part of the experimental discipline.

\section{Chromatic color test already visible in DES}
The opacity law predicts an intrinsic redward trend approximately
\begin{equation}
 \frac{dc_{\rm prop}}{dz}\simeq+0.019\ {\rm mag}/z,
 \label{eq:colorpred}
\end{equation}
with representative $E(B-V)\simeq0.007,0.0125,0.020$ at $z=0.25,0.5,1$ under a simple monochromatic mapping. On 1,635 DES-subset SNe the raw SALT color slope is instead
\begin{equation}
 \frac{dc_{\rm raw}}{dz}=-0.0963\pm0.0140\ {\rm mag}/z.
 \label{eq:colorraw}
\end{equation}
The opposite sign is not presently a clean propagation falsification because the flux-limited DES sample becomes strongly bluer with redshift under its own selection pipeline; the predicted signal is roughly five times smaller than the raw selection trend. The immediate inexpensive comparison is data minus the published FLRW+selection simulation as a function of redshift. A full alternative-reference BBC simulation remains the stronger end-to-end test. Version 8 therefore reports Eq.~\eqref{eq:colorraw} rather than describing color as an unexamined future observable.

\section{Tolman surface brightness and radiance}
\label{sec:tolman}
With static Euclidean geometry, a strictly forward coherent channel, and no transverse magnification from the tiny transient group-index perturbation,
\begin{equation}
 D_A=D=\frac{c}{K}\ln(1+z).
 \label{eq:DA}
\end{equation}
The energy-redshift and arrival-rate factors give, before non-forward attenuation,
\begin{equation}
 D_{L,0}=(1+z)D,
 \qquad
 \frac{D_{L,0}}{D_A}=1+z,
 \label{eq:duality}
\end{equation}
not the metric Etherington ratio $(1+z)^2$. Including Eq.~\eqref{eq:tau},
\begin{equation}
 I_{\rm obs}\propto(1+z)^{-2}e^{-\tau_\lambda(z)}.
 \label{eq:tolman}
\end{equation}
Define the effective exponent by $I\propto(1+z)^{-n_{\rm eff}}$,
\begin{equation}
 n_{\rm eff}(z,\lambda)=2+\frac{\tau_\lambda(z)}{\ln(1+z)}.
 \label{eq:neff}
\end{equation}
Near the Lubin--Sandage cluster mean $z\simeq0.86$, the model predicts approximately $n_R\simeq2.18$ and $n_I\simeq2.16$.

Lubin and Sandage analyzed 34 early-type galaxies in three clusters near $z=0.8$--0.9 and reported $n_R=2.59\pm0.17$ and $n_I=3.37\pm0.13$ in their $q_0=1/2$ distance frame \citep{LubinSandageIII,LubinSandageIV}. They also published the inferred exponents for other $q_0$ choices, directly exposing the distance-frame sensitivity of the placement on the local surface-brightness--linear-radius relation.

Before retrieving a full per-galaxy static reduction, we preregistered an analytic frame-shift diagnostic using that published $q_0$ grid. The measured response of the inferred exponent to $\log D_A$ is approximately $-3.41$/dex in $R$ and $-3.17$/dex in $I$. Extrapolating to Eq.~\eqref{eq:DA} moves the nominal values to
\begin{equation}
 n_R\simeq1.45\pm0.17,
 \qquad
 n_I\simeq2.30\pm0.13.
 \label{eq:frameShift}
\end{equation}
Thus the often-quoted $I$-band crisis is highly frame dependent, but the problem does not disappear: it migrates to $R$. More importantly, the two bands move similarly under distance reassignment, leaving the approximately frame-invariant spread
\begin{equation}
 (n_I-n_R)_{\rm L\&S}=0.78\pm0.21
 \label{eq:spread}
\end{equation}
against the model prediction of approximately $-0.02$. This is a 3.7$\sigma$ discrepancy before a modern systematic reanalysis.

The analytic extrapolation extends about 3.5 times beyond the $D_A$ range over which the published $q_0$ grid calibrates the response, so Eq.~\eqref{eq:frameShift} is not a resolution. The full static-frame reduction is preregistered: re-place the observed angular Petrosian radii using Eq.~\eqref{eq:DA}, recompute transfer/K-corrections consistently, and report the $I$ improvement, $R$ migration, and band spread together.

Version 8 also corrects an overrestriction in the original preregistration. Static geometry does not logically imply zero luminosity evolution because Eq.~\eqref{eq:lookback} is finite. Two branches are therefore preserved before data reduction: (Z) the originally registered zero-evolution test, and (E) stellar-population evolution computed independently at the static lookback time, with no evolution amplitude fitted to the Tolman residuals. If independently constrained evolution cannot reproduce the sign and magnitude of the band spread, it cannot be invoked post hoc as a rescue.

The same duality difference creates a radiance test. A localized thermal source whose spectral coordinates are shifted to $T/(1+z)$ is, before opacity, brighter by $(1+z)^2$ than a true Planck surface at that shifted temperature. Expanding-photosphere measurements of Type II supernovae at moderate redshift therefore provide a separate distance/radiance discriminator.

\section{Global evolution is a real constraint, not a forbidden parameter}
The observation that galaxy populations, star-formation activity, and quasar populations vary strongly with redshift is not erased by static geometry. Large surveys show pronounced evolution in the cosmic star-formation-rate density \citep{MadauDickinson2014}. Version 8 therefore does not claim a strictly stationary eternal matter distribution.

The adopted branch is: spatial geometry is static, while source populations may evolve with emission epoch. Their evolution must be described against the static lookback law Eq.~\eqref{eq:lookback}, not fitted independently to each cosmological observable. This creates a new class of tests: stellar ages, fundamental-plane evolution, star-formation history, and quasar demographics can be compared under the static and FLRW lookback functions. At $z\sim0.9$ the difference is modest enough that current stellar ages may not discriminate sharply; at higher redshift the histories diverge more strongly. No global population-history model is claimed here.

Quasar variability now supplies a response-bandwidth constraint rather than the earlier proposed ``steady-source under-dilation'' signature, which is withdrawn. Brewer et al. (2025) find a cosmological dependence $(1+z)^n$ with $n=1.14\pm0.34$, consistent with unity and robust to their alternative modeling assumptions \citep{BrewerLewisLi2025}; Lewis \& Brewer (2023) found a similar result \citep{LewisBrewer2023}. A quasar's variability is an AC perturbation $\delta f_{\hat{k}}(t)$ about a persistent DC source, so the direction-resolved reactive response may track the fluctuations while the DC component remains saturated. Existing decade-baseline measurements therefore constrain the bandwidth, saturation, and relaxation of $\mathcal{R}$; if interpreted as one persistent memory cell they favor relaxation/memory scales extending to at least decades. Retiring generic quasar under-dilation also removes it as a wake-versus-Gordon discriminator: the surviving distinction is geometric, principally $D_A=D$ and the $(1+z)$ distance-duality ratio rather than FLRW's $(1+z)^2$.

\section{Forward and non-forward energy closure}
\label{sec:energy}
The non-forward ledger accounts explicitly for 16\% rest-$B$ beam depletion at $z=1$, but the larger coherent energy transfer in Eq.~\eqref{eq:fwdloss} must be included in the same accounting. At $z=1$, half the emitted photon energy has left the observed photon through the forward frequency shift. The theory therefore owes the combined ledger
\begin{equation}
 E_{\rm emitted}=E_{\rm forward,obs}+E_{\rm nonforward}+E_{\rm bath}+E_{\rm stored},
 \label{eq:ledger}
\end{equation}
with the stored wake energy remaining bounded under steady illumination.

The microwave bath itself cannot be treated as a stationary fixed point if Eq.~\eqref{eq:drag} is universal. For $T_0=2.7255$ K and $u_{\rm CMB}\simeq4.17\times10^{-13}$ erg cm$^{-3}$, its direct frequency-drift term has magnitude
\begin{equation}
 K u_{\rm CMB}\simeq9.5\times10^{-31}\ {\rm erg\,cm^{-3}\,s^{-1}},
 \label{eq:cmbdriftenergy}
\end{equation}
whereas taking the optical+IR extragalactic background at roughly 50 nW m$^{-2}$ sr$^{-1}$ gives $u_{\rm opt}\sim2.1\times10^{-14}$ erg cm$^{-3}$ and
\begin{equation}
 \dot u_{\rm opt\rightarrow bath}\sim u_{\rm opt}K\sim4.8\times10^{-32}\ {\rm erg\,cm^{-3}\,s^{-1}}.
 \label{eq:optinject}
\end{equation}
The bath-drift term is therefore about twenty times larger than the identified optical injection. This is not a steady-maintenance problem if the bath itself evolves. Eq.~\eqref{eq:drag} together with the static lookback law predicts
\begin{equation}
 \boxed{T_{\rm CMB}(z)=T_0(1+z)}.
 \label{eq:tcmbz}
\end{equation}
At $z=2.41837$, Eq.~\eqref{eq:tcmbz} gives 9.315 K; CO rotational excitation gives $9.15\pm0.72$ K \citep{Srianand2008}. A broader CO+SZ analysis finds $T_{\rm CMB}(z)=(2.725\pm0.002)(1+z)^{1-\beta}$ K with $\beta=-0.007\pm0.027$ \citep{Noterdaeme2011}. Thus the measured $T(z)$ relation is a zero-additional-parameter \emph{consistency success} of the drift law, but not a discriminator: FLRW predicts the same scaling.

Frequency drift alone does not preserve the absolute normalization of a blackbody in a static volume. Because a Planck photon number density scales as $T^3$, a self-similar drifting Planck bath requires a companion number-relaxation law
\begin{equation}
 \frac{d\ln n_\gamma}{dt}=-3K
 \label{eq:numberrelax}
\end{equation}
(or an equivalent kinetic term) while the frequency coordinates drift at $-K$. The required number-changing process is \textbf{not identified in this paper}. It must be broadband, shape preserving to FIRAS precision, and must not introduce observable attenuation or chromaticity into directed beams. A proposed low-frequency cutoff that simply exempts microwave photons is also disfavored independently: a thermal factor $1-e^{-h\nu/kT}$ would reduce the drag at the 1.420 GHz H I line to about 2.5\% of its optical value, whereas ensemble 21-cm/UV absorber comparisons constrain differential radio/optical redshift behavior at the $\sim10^{-5}$ level (with individual astrophysical velocity offsets of order km s$^{-1}$) \citep{Tzanavaris2007}. Such a cutoff is therefore not used.

The combined bath picture generates explicit falsification gates:
\begin{enumerate}[label=(\roman*)]
\item the non-forward differential kernel $d^2\sigma/(d\Omega\,dE')$ must not create redshift-correlated optical PSF wings or an optical/IR background excess;
\item the direction-resolved reactive memory must saturate and relax rather than grow without bound;
\item the number-changing/thermalization kinetics must preserve a Planck spectrum and its absolute normalization to FIRAS precision while allowing Eq.~\eqref{eq:tcmbz};
\item the same bath dynamics must reproduce, not merely coexist with, the observed CMB acoustic/polarization structure, lensing, and its relation to the BAO scale.
\end{enumerate}

COBE/FIRAS constrains RMS deviations from a blackbody to below roughly 50 ppm of the peak, with $|y|<15\times10^{-6}$ and $|\mu|<9\times10^{-5}$ \citep{Fixsen1996}. A previous 25 Myr estimate based only on thermalizing injected optical energy remains an order-of-magnitude sub-gate for the injected component, but it does not close the larger number-normalization problem in Eq.~\eqref{eq:numberrelax}. A microscopic model must compute the full collision operator and test FIRAS directly.

The CMB is therefore treated here as a \emph{drifting self-similar bath}, not as a stationary equilibrium field and not as a 3000 K last-scattering photosphere propagated through the localized-source surface-brightness relation. Its measured temperature-redshift law is consistent with the drift; its present-day Planck normalization and anisotropy structure remain active theoretical barriers. Planck's final anisotropy data are well described by six-parameter base \LCDM \citep{Planck2018}. Until an alternative bath dynamics produces comparable spectral and angular structure, the present model is an incomplete cosmology even if its SN transfer law survives.

\section{Purpose-built photon-counting experiment}
The chromatic prediction can be tested more directly than conventional classification spectroscopy permits. Once a supernova is discovered and its redshift measured, a purpose-built telescope can monitor selected observer-frame passbands corresponding to fixed rest-frame spectral windows.

A practical instrument would use a sufficiently large collecting mirror, dichroic splitting or narrow-band filters, and simultaneous APD/SPAD photon-counting channels. Red-enhanced silicon detectors cover much of the optical range; moderate-redshift monitoring of red rest-frame windows may require a near-IR counting channel. For a target at known $z$,
\begin{equation}
 \lambda_{\rm filt}\simeq(1+z)\lambda_{\rm rest}.
\end{equation}
The experiment should avoid claims about isolated Ni/Co/Fe lines in broad SN spectra. Instead it can follow integrated phase-matched windows such as the iron-group-dominated 320--340 nm region and a 620--680 nm red reference, plus an $i/z$/near-IR channel containing the secondary maximum.

The core observable is a simultaneous photon-count ratio
\begin{equation}
 R_{12}(t,z)=\frac{N_1(t,z)}{N_2(t,z)}.
\end{equation}
After throughput calibration with field standards and an internal source, the propagation law predicts
\begin{equation}
 \frac{R_{12}(z)}{R_{12}(0)}=
 \exp\left\{-\tauinf\left[\sqrt{\frac{440}{\lambda_1}}-\sqrt{\frac{440}{\lambda_2}}\right]
 \left[1-(1+z)^{-1/2}\right]\right\},
 \label{eq:instrument}
\end{equation}
with no new opacity amplitude fitted to this experiment. The cosmological component must be approximately phase-flat at matched rest wavelengths. A sufficiently precise null therefore kills the non-forward $\lambda^{-1/2}$ channel. High cadence simultaneously improves the temporal-clock test that DES could not decide.

Before observing, the exact rest windows, filter widths, detector efficiencies, cadence, calibration-star protocol, estimator, redshift bins, and kill threshold should be preregistered. The instrument is useful precisely because it can return a meaningful null.

\section{Standalone falsification ledger}
\small
\begin{longtable}{p{0.20\linewidth}p{0.25\linewidth}p{0.47\linewidth}}
\caption{Current status of the principal claims. ``Consistent'' is not evidence favoring the model when FLRW predicts the same observable.}\label{tab:ledger}\\
\toprule
Observable/test & Current result & Status and failure condition\\
\midrule
\endfirsthead
\toprule
Observable/test & Current result & Status and failure condition\\
\midrule
\endhead
DES Hubble diagram & $\Delta\chi^2=+1.1$ vs flat \LCDM & Near degeneracy on released SN distances; SN-only result, not a global cosmological fit.\\
Achromatic spectral redshift & common $z$ across spectral features & Required by coherent face; classical strongly chromatic loss mechanisms excluded. Microscopic vertex rate $K$ still owed.\\
DES bulk duration & $b=1.003\pm0.005\pm0.010$ & FLRW predicts $b=1$; static model has $b=\kreact$ with $\kreact$ empirically constrained $\simeq1$. Equality to unity is not yet derived microscopically.\\
Spectral aging & aging $\propto1/(1+z)$ & Independent SN spectral clock; consistent with universal remap and FLRW.\\
Filter/Wien artifact & wake off: $|\Delta s|\le5\times10^{-4}$ & \textbf{Null result: proposed explanation falsified}. Chromatic erosion cannot manufacture $b=1$.\\
T4 secondary maximum & $b_2=0.52\pm0.78$, $N=120$ & \textbf{Inconclusive}; contains 0 and 1. DES cadence/redshift range underpowered.\\
Raw DES color trend & $-0.0963\pm0.0140$ mag/$z$ vs intrinsic propagation $+0.019$ & Opposite raw sign, but selection trend is larger; compare data-minus-published selection simulation now, then matched alternative-reference BBC.\\
Chromatic phase-flat ratio & $Q=0.979/0.964/0.944$ at $z=.25/.5/1$ & Open direct discriminator; amplitude fixed by Hubble fit. A precise phase-flat null kills $\lambda^{-1/2}$ channel.\\
Tolman distance frame & static extrapolation gives $n_I\simeq2.30$, $n_R\simeq1.45$ & $I$ crisis is frame fragile; $R$ tension appears. Full per-galaxy reduction preregistered.\\
Tolman band spread & $0.78\pm0.21$ observed vs $-0.02$ predicted & \textbf{Strongest current empirical threat} ($\sim3.7\sigma$ before modern systematics/evolution calculation).\\
Population evolution & static lookback $K^{-1}\ln(1+z)$ & V8 adopts evolving contents; must reproduce observed stellar/SFR/quasar evolution without fitting each observable independently.\\
Forward energy transfer & $z/(1+z)$; 50\% at $z=1$ & Major open ledger. Energy destination, bath number relaxation, and full kinetic closure must be derived.\\
Non-forward beam depletion & 10.3\% at $z=.5$, 16.0\% at $z=1$ & Differential angular/spectral kernel must satisfy halo and EBL limits.\\
CMB temperature drift & $T(z)=T_0(1+z)$; 9.315 K predicted vs $9.15\pm0.72$ K at $z=2.418$ & \textbf{Consistency success shared with FLRW}, not discrimination. Stationary 2.725 K bath at all lookbacks is excluded.\\
FIRAS normalization & $<50$ ppm RMS; $|\mu|<9\times10^{-5}$ & \textbf{Active theoretical barrier}. Drifting Planck bath requires shape-preserving photon-number relaxation at rate $3K$; mechanism not identified.\\
CMB anisotropy/BAO & not computed & \textbf{Major unresolved cross-sector test}. Bath must reproduce acoustic/polarization/lensing structure and its relation to BAO.\\
Two-face microphysics & $\kreact\simeq1$ empirically constrained by DES & Derive coherent inelastic drift, direction-resolved reactive response, bandwidth/relaxation, cell selection, and coefficient equality from one vacuum action; otherwise model remains effective phenomenology.\\
\bottomrule
\end{longtable}
\normalsize

\section{What Version 8 establishes and what it does not}
Version 8 does not establish a static universe. It establishes something narrower but nontrivial: a specific static-geometry supernova model can be written without the two immediate defects that kill classical tired light---chromatic redshift and $b=0$ time behavior---and can be confronted with modern DES data at essentially the same SN Hubble-diagram fit quality as flat \LCDM. The observational law is no longer an arbitrary curve: it comes with explicit timing, chromatic, angular-distance, radiance, and energy consequences.

The adversarial process also removed attractive but false support. Passband/Wien erosion does not create time dilation. T4 did not confirm the universal remap; it was underpowered. The Tolman band spread is adverse. The raw DES color trend is opposite the intrinsic chromatic prediction but presently selection dominated. The forward energy loss is larger than the non-forward loss and creates a serious CMB number-normalization, thermalization, and anisotropy debt. These are not footnotes to be edited away; they define the remaining research program.

The two-face sourcing model is the principal theoretical refinement. A universal coherent inelastic face supplies the frequency drift, while a direction-resolved phase-space susceptibility can remap intervals between independent photons without requiring one scalar wake to remember every transient crossing a spatial point. This resolves the earlier sourcing and cross-talk problems at the level of an effective transport model. It does not yet solve the microscopic problem: the same vacuum dynamics must produce the required universal $K$, the empirical reactive coefficient $\kreact\simeq1$, directional/polarization cell selection, appropriate bandwidth and relaxation, the $3K$ bath number-relaxation companion, and the correct angular/spectral redistribution. If no such transport response can be derived, the model fails theoretically even if the one-dimensional SN fit remains good.

Parsimony is therefore not claimed globally. Base \LCDM has its own parameter and modeling structure but is already tested across the CMB, BAO, lensing, growth, and nucleosynthesis. Removing dark-energy acceleration from the SN interpretation is a parsimony gain only if the replacement propagation physics ultimately closes with comparable or lower total structure. Version 8 leaves that comparison open.

\section{Conclusion}
The DES-SN5YR Hubble diagram alone remains nearly indifferent between one-parameter flat \LCDM and the one-parameter static two-channel distance law, with $\Delta\chi^2=+1.1$. That is the empirical wedge, not the conclusion of the cosmology.

The model now has a standalone causal timing formulation. The coherent inelastic face supplies the universal frequency drift of Eq.~\eqref{eq:drag}; temporal structure in a direction-resolved phase-space cell generates the separate group-delay response of Eq.~\eqref{eq:ng}, shared by later independent photons in that cell. The resulting temporal exponent is $b=\kreact$, and DES empirically constrains $\kreact\simeq1$; FLRW predicts unity without this extra response coefficient. Deriving the equality of the two rates from microscopic dynamics is therefore a central theory-side target rather than a hidden assumption.

The non-forward $E^{1/2}$ channel supplies the fitted Hubble curvature and predicts a fixed, weak chromatic erosion. Its alternative use as an explanation for time dilation was explicitly tested and falsified. The preregistered T4 species-clock measurement was honestly inconclusive. The Tolman $R-I$ spread remains the most immediate adverse optical test. The forward channel transfers 50\% of photon energy by $z=1$, making bath number normalization, thermalization, and CMB structure unavoidable parts of the same model rather than optional future topics.

The next useful work is therefore sharply defined: execute the full preregistered Tolman re-reduction; perform the data-minus-DES-selection color comparison and ultimately the matched alternative-reference BBC simulation; derive the two-face vertex and its energy redistribution from the vacuum action; test the drifting bath and its number-relaxation kinetics against FIRAS and CMB anisotropies; and build the high-cadence multi-channel photon-counting experiment whose null can directly eliminate the chromatic channel.

A theory that survives only because every contradiction is reclassified is scientifically empty. Version 8 is organized to prevent that outcome: the model either closes these gates or it fails.

\section*{Data, code, and version provenance}
The public DES-SN5YR light curves and v1.2 distance products are the observational basis of this work. The canonical code repository, effective-field notes, transfer-integral producers, preregistrations, T4 execution report, Tolman frame-shift analysis, forward-energy ledger, and instrument note are maintained at \href{https://github.com/tracyphasespace/Static-Universe-DES-SN5YR}{github.com/tracyphasespace/Static-Universe-DES-SN5YR}.

Version 8 is archived at DOI \href{https://doi.org/\veightdoi}{\veightdoi}. Prior records are intentionally preserved rather than overwritten: Version 7 at \href{https://doi.org/\vsevendoi}{\vsevendoi}, Draft 6 at \href{https://doi.org/\draftsixdoi}{\draftsixdoi}, and Draft 5 at \href{https://doi.org/\draftfivedoi}{\draftfivedoi}. The independent BBC/model-discrimination methods analysis is archived at DOI \href{https://doi.org/\methodsdoi}{\methodsdoi}. These records document how adverse review, null tests, and implementation corrections changed the model.

\section*{Acknowledgments}
The author thanks the Dark Energy Survey collaboration for making the light curves, simulations, covariance products, and analysis documentation public. AI-assisted tools were used for adversarial review, code checking, literature cross-checking, and manuscript drafting. Failed mechanisms, preregistration amendments, implementation defects, inconclusive measurements, and adverse tests are retained in the public record. The author is responsible for the hypotheses, calculations, interpretation, and final text.

\begin{thebibliography}{99}
\bibitem[Bamber et al.(1999)]{Bamber1999} Bamber, C., Boege, S.~J., Koffas, T., et al. 1999, Phys. Rev. D, 60, 092004.
\bibitem[Blondin et al.(2008)]{Blondin2008} Blondin, S., Davis, T.~M., Krisciunas, K., et al. 2008, ApJ, 682, 724, doi:10.1086/589568.
\bibitem[Brewer et al.(2025)]{BrewerLewisLi2025} Brewer, B.~J., Lewis, G.~F., \& Li, Y. 2025, MNRAS, 537, 809, doi:10.1093/mnras/staf044.
\bibitem[Burke et al.(1997)]{Burke1997} Burke, D.~L., Field, R.~C., Horton-Smith, G., et al. 1997, Phys. Rev. Lett., 79, 1626.
\bibitem[Camilleri et al.(2024)]{Camilleri2024} Camilleri, R., Davis, T.~M., Vincenzi, M., et al. 2024, MNRAS, 533, 2615.
\bibitem[Srianand et al.(2008)]{Srianand2008} Srianand, R., Noterdaeme, P., Ledoux, C., \& Petitjean, P. 2008, A\&A, 482, L39, doi:10.1051/0004-6361:200809727.
\bibitem[Noterdaeme et al.(2011)]{Noterdaeme2011} Noterdaeme, P., Petitjean, P., Srianand, R., Ledoux, C., \& L\'opez, S. 2011, A\&A, 526, L7, doi:10.1051/0004-6361/201016140.
\bibitem[Tzanavaris et al.(2007)]{Tzanavaris2007} Tzanavaris, P., Murphy, M.~T., Webb, J.~K., Flambaum, V.~V., \& Curran, S.~J. 2007, MNRAS, 374, 634, doi:10.1111/j.1365-2966.2006.11171.x.
\bibitem[Fixsen et al.(1996)]{Fixsen1996} Fixsen, D.~J., Cheng, E.~S., Gales, J.~M., Mather, J.~C., Shafer, R.~A., \& Wright, E.~L. 1996, ApJ, 473, 576.
\bibitem[Hawkins(2010)]{Hawkins2010} Hawkins, M.~R.~S. 2010, MNRAS, 405, 1940.
\bibitem[Hsiao et al.(2007)]{Hsiao2007} Hsiao, E.~Y., Conley, A., Howell, D.~A., et al. 2007, ApJ, 663, 1187.
\bibitem[Jack et al.(2015)]{Jack2015} Jack, D., Baron, E., \& Hauschildt, P.~H. 2015, MNRAS, 449, 3581, doi:10.1093/mnras/stv474.
\bibitem[Lewis \& Brewer(2023)]{LewisBrewer2023} Lewis, G.~F., \& Brewer, B.~J. 2023, Nature Astronomy, 7, 1265, doi:10.1038/s41550-023-02029-2.
\bibitem[Lubin \& Sandage(2001a)]{LubinSandageIII} Lubin, L.~M., \& Sandage, A. 2001a, AJ, 122, 1071, doi:10.1086/322133.
\bibitem[Lubin \& Sandage(2001b)]{LubinSandageIV} Lubin, L.~M., \& Sandage, A. 2001b, AJ, 122, 1084, doi:10.1086/322134.
\bibitem[Madau \& Dickinson(2014)]{MadauDickinson2014} Madau, P., \& Dickinson, M. 2014, ARA\&A, 52, 415, doi:10.1146/annurev-astro-081811-125615.
\bibitem[Planck Collaboration(2020)]{Planck2018} Planck Collaboration VI. 2020, A\&A, 641, A6, arXiv:1807.06209.
\bibitem[S\'anchez et al.(2024)]{Sanchez2024} S\'anchez, B.~O., Brout, D., Vincenzi, M., et al. 2024, ApJ, 975, 5.
\bibitem[Vincenzi et al.(2024)]{Vincenzi2024} Vincenzi, M., Brout, D., Armstrong, P., et al. 2024, ApJ, 975, 86.
\bibitem[White et al.(2024)]{White2024} White, R.~M.~T., Davis, T.~M., Lewis, G.~F., et al. 2024, MNRAS, 533, 3365.
\bibitem[Zwicky(1929)]{Zwicky1929} Zwicky, F. 1929, Proc. Natl. Acad. Sci. USA, 15, 773.
\end{thebibliography}
\end{document}
